\documentclass[10pt,a4paper]{amsart}
\usepackage[margin=19mm]{geometry}
\usepackage{amsmath,amssymb,mathtools}
\usepackage[scr=rsfs,cal=euler]{mathalfa}
\usepackage{xfrac}
\usepackage[unicode,hidelinks,bookmarks=false]{hyperref}
\newcommand{\ie}{i.e.\ }
\newcommand{\cf}{cf.\ }
\newcommand{\resp}{respectively\ }
\newcommand{\Subsection}{\subsection}
\providecommand{\nobreakdash}{\nobreak\hbox{-}}
\setlength{\emergencystretch}{2em}
\multlinegap0pt
\usepackage{amssymb}
\usepackage{tikz}
\newtheorem{lemma}{Lemma}
\title{Backward Arcs in Hamilton Oriented Cycles and Paths in Directed Graphs with Independence Number Two}
\author{S. Gerke, Q. Guo, G. Gutin, Y. Hao, W. Veeranonchai, A. Yeo}
\date{Source: arXiv:2603.22993v1 (CC BY 4.0)}
\begin{document}
\maketitle

\noindent\emph{Source and licence.} Backward Arcs in Hamilton Oriented Cycles and Paths in Directed Graphs with Independence Number Two. Version arXiv:2603.22993v1 (CC BY 4.0), distributed by arXiv under the Creative Commons Attribution 4.0 International licence, http://creativecommons.org/licenses/by/4.0/, as recorded in the arXiv licence metadata for this exact version and shown on the abstract page https://arxiv.org/abs/2603.22993v1. Full licence notice: \texttt{licenses/arxiv-2603.22993-CC-BY-4.0.txt}.\par\smallskip

\noindent\textit{Editorial scope.} Counted unit: the article's lemma XYpathSemiC (source0.tex lines 592--599). The article's complete original proof of it is delivered with it (source0.tex lines 601--660, one \texttt{proof} environment of 60 lines). No mathematics is written, repaired or re-derived: the statement and the proof are the article's own lines, in the article's own notation, and no retained line is substituted or re-flowed.  No further source-local context is needed: every cross-reference of every retained line resolves inside the two delivered spans.  Nothing was omitted from any retained span, so the package carries no omission record.  No retained line contains a citation, so the wrapper carries no bibliography and no bibliography entry.  No retained line contains a figure environment, an image-inclusion command or drawing code, so no image is delivered and the package declares no asset dependency.  Editorial formatting only, in addition: an AMS \texttt{amsart} preamble that restates the article's own declarations and macros used by the retained lines, namely the environments \texttt{lemma} on separate counters, together with the standard packages those lines need.  The retained environments print in this document as Lemma 1 in order of appearance, whereas the article prints the same items with the numbering of its own full document.  The complete native source of the article as distributed in the arXiv e-print for this exact version is bundled unchanged as \texttt{sources/197046/source0.tex}.
\begin{lemma} \label{lem: XYpathSemiC}
Let $D$ be a semicomplete digraph and let $x,y \in V(D)$ be arbitrary distinct vertices. Then 
\begin{description}
\item[(a)] $b(x,y)\le2$.

\item[(b)] $b(x,y)\le 1$ or $b(y,x)\le 1$.
\end{description}
\end{lemma}

\begin{proof}
We first prove part~(a).
If $|V(D)| \le 3$ then clearly the lemma holds, so assume that  $|V(D)| \geq 4$.
Let $D' = D - y$ and note that adding at most one arc to $D'$, we can make $D'$ strong.  Therefore, there exists a Hamilton orcycle in $D'$ with 
at most one backward arc. Deleting the arc between $x$'s predecessor, $x^-$, and $x$ and adding the arc between $x^-$ and $y$ instead gives us the desired orpath in $D$.
This completes the proof of part~(a).

Now we prove part~(b).
First, suppose that $D$ is strong. Let $C$ be a Hamilton dicycle in $D$ and let $x,y \in V(D)$ be arbitrary distinct vertices.

We denote $x^+,y^+$ the successors of $x,y$ on $C$, and $x^{++},y^{++}$ the successors of $x^+,y^+$ on $C$.

\vspace{0.5em}
\textbf{Case 1:} $x^+=y$ or $x=y^+$.

There is a Hamilton dipath between $x$ and $y$ in $D$.

\vspace{0.5em}
\textbf{Case 2:} $x^{++}=y$ or $x=y^{++}$, and $x^+\ne y$ and $x\ne y^+$.

If $x^{++}=y$ and $x=y^{++}$, then $|V(D)|=4$ and the length of any Hamilton orpath between $x$ and $y$ is $3$, so this orpath or its reverse is the desired orpath.

If not, without loss of generality, $x^{++}=y$ and $x\ne y^{++}$, and then  $y$'s predecessor on $C$, $y^-$ is $x^+$. Then one of the following two Hamilton orpaths is the desired orpath, depending on the arc between $x^+$ and $y^+$.

\[
P_1=x x^+ C[y^+, x^-] y,~P_2=y y^+ x^+ C[y^{++}, x]
\]

\vspace{0.5em}
\textbf{Case 3:} $x^{++}\ne y$, $x\ne y^{++}$, $x^+\ne y$, and $x\ne y^+$.

Then one of the following two Hamilton orpaths is the desired orpath, depending on the arc between $x^+$ and $y^+$.
\[
P_1=x x^+ C[y^+, x^-] C[x^{++}, y], ~P_2=y y^+ C[x^+, y^-]C[y^{++}, x]
.\]
This completes the case when $D$ is strong.

Now consider the case when that $D$ is not strong.
Decompose $D$ into strong components $V_1,\dots,V_l$ with all arcs pointing from $V_i$ to $V_j$ whenever $i<j$ ($l>1$). All strong semicomplete digraphs contain a Hamilton dicycle, and hence contain a Hamilton dipath starting or ending at any vertex (not necessarily at the same time). We denote $P_i$ to be an arbitrary Hamilton dipath in $V_i$, and for a vertex $v$, we denote $S_v$ ($E_v$, respectively) to be a Hamilton dipath in the strong component containing $v$ starting at $v$ (ending at $v$, respectively).

Let $x \in V_a$ and $y \in V_b$.
Without loss of generality, we may assume that $a \le b$.

\textbf{Case 1:} $a<b$.
In this case, the desired Hamilton $(x,y)$-orpath is the following:
\[
S_xP_{a+1}\dots P_{b-1}P_{b+1}\dots P_lP_1 \dots P_{a-1}E_y
\]
The only backward arc is the one going from $P_l$ to $P_1$.

\textbf{Case 2:} $a=b$. %By the strong case, and 
Without loss of generality, we may assume that $b(x,y)\le 1$ in the strong component $V_a$.
Let $u_1u_2\dots u_s$, with $u_1=x, u_s=y$, be the Hamilton $(x,y)$-orpath with at most one backward arc in $V_a$.
Say the backward arc (if it exists) is between $u_i$ and $u_{i+1}$ for $i \in \{1,2,\dots s-1 \}$. 
Then the desired Hamilton $(x,y)$-orpath in $D$ is the following:
\[
xu_2\dots u_{i}P_{a+1}\dots P_lP_1 \dots P_{a-1}u_{i+1}\dots u_{s-1}y.
\]
Again, the only backward arc is the one going from $P_l$ to $P_1$.
\end{proof}   

\end{document}
